%!TEX root = ../../thesis-trier.tex

Here, we investigate the constrained synchronization problem
for binary SI-automata as input.
\Cref{constrsync:input:SI:sec:monoid} collects results
for constraint languages that are monoids.
\Cref{constrsync:input:SI:sec:bounded} collects results
for bounded constraint languages.
Lastly, \cref{constrsync:input:SI:sec:constraintaut}
shows that for constraint PDFAs with up to three states
and binary SI-automata as input 
the constrained synchronization problem is in \NL.
 
\begin{center} 
	\fbox{
		\begin{minipage}{30em}
			In the remainder of this section we denote by $\mathcal A_t$ the synchronizing
			binary SI-automaton from~\cref{prelim:fig:binarySI-sync-trap} with a trap state
			and by $\mathcal A_c$ the other automaton from~\cref{prelim:fig:binarySI-sync-SCC}
			where $b$ permutes the states,~$a$ is a simple idempotent
			and $s,t \in Q$ and $0 < d < n$
			are the parameters as in \cref{prop:structure_simp_idem_binary} w.r.t.~$\mathcal A_t$ and $\mathcal A_c$
			for $n \ge 4$.
		\end{minipage}
	}
\end{center}
\label{def:AtAc}

\subsection{Monoid Constraint Languages} 
\label{constrsync:input:SI:sec:monoid}

Here, we collect various results on monoid languages.
For example, \cref{constrsync:subclasses:SI:prop:uAc} characterizes those words $u \in b^*ab^*$ 
with precisely one $a$ such that~$u^*$ contains a synchronizing word for~$\mathcal A_c$.

% für m > 2 sogar mengen die von a und bestimmten wörtern permutiert.
\begin{proposition} \label{constrsync:subclasses:prop:SI:permutesset-new}
	Suppose $L \subseteq ( (ba^*)^{m - 1}b )^*$
	for a divisor $m$ of~$n$ (resp.~$n- 1$)
	with $m > 1$.
	Then there exists a proper subset of states of size $n / m$ (resp.~$(n-1)/m$)
	permuted by each word from~$L$ in $\mathcal A_c$ (resp.~$\mathcal A_t$).
\end{proposition}
\begin{proof}
	Recall that $s \in Q$ denotes the unique state with $\delta(s, a) \ne s$.
	Set $S = \{ \delta(s, b^{mi}) \mid i \ge 0 \}$.
	Then $|S| = n / m$ (resp.~$(n-1)/m$) in $\mathcal A_c$ (resp.~$\mathcal A_t$).
	We have $s \notin \delta(S, b^i)$ for all $1 \le i \le m - 1$.
	Thus $\delta(q, a) = q$ for all $q \in \delta(S, b^i)$ with $1 \le i \le m - 1$.
	This gives $\delta(S, u) = \delta(S, u^{|u|_b}) = S$
	for all $u \in L$.	
\end{proof}

\begin{corollary} \label{constrsync:subclasses:cor:SI:permutesset-new}
	For $\mathcal A_c$ and $n$ even (resp.~$\mathcal A_t$ and $n$ odd) %and $n \ge 4$
	there exists a proper subset of states of size $n / 2$ (resp.~$(n-1) / 2$) permuted by each word in $(ba^*b)^*$.
\end{corollary}

The next statement generalizes 
\crefitem{prelim:prop:binary-SI-sync-words}{prelim:prop:binary-SI-sync-words:1}.

\begin{proposition} \label{constrsync:subclasses:SI:prop:u-coprime-bs}
	Let $u \in \{a,b\}^*$ with $|u|_a > 0$ and $|u|_b > 0$
	and suppose $|u|_b$ and $n - 1$ are coprime. Then
	$\Syn(\mathcal A_t) \cap u^* \ne \emptyset$.
\end{proposition}
\begin{proof}
	As $|u|_b$ and $n - 1$ are coprime,
	there exists $i \ge 1$ such that $i\cdot |u|_b \equiv 1 \pmod{n - 1}$
	and so~$b^{i \cdot |u|_b}$ and $b$ induce
	the same permutation on $Q$.
	Hence, for each $q \in Q$, we have
	$\delta(q, u^i) \in \{t, \delta(q, b)\}$.
	As $|u|_a > 0$ there exists $q \in Q \setminus \{t \}$
	such that $\delta(q, u) = t$.
	So, we obtain that for each  $p \in Q \setminus \{t\}$ there exists $0 \le j < n - 1$
	such that $\delta(p, u^j) \in  \{t,q \}$,
	which gives $\delta(p, u^{j+1}) = t$.
	Hence $\delta(p, u^{n-1}) = t$ for all $p \in Q$
	and so $u^{n-1} \in \Syn(\mathcal A_t)$.    
\end{proof}

\begin{proposition}
	\label{constrsync:subclasses:SI:prop:uAc}
	Let $u \in b^*ab^*$.
	%
	% $i\cdot |u|_b + d \equiv 0 \pmod{n}$ for some $i \ge 1$.
	% permutiert dann s,s.u, ... s.u^i, weil s.u^i = s mit der wahl
	Then $\Syn(\mathcal A_c) \cap u^* \ne \emptyset$
	if and only if $d + |u|_b \equiv 0 \pmod{n}$.
\end{proposition}
\begin{proof}
	We have $\fact{u^*} = \fact{(ab^{|u|_b})^*}$.
	So, by~\cref{constr_sync:cor:factors_sets_included},
	we can assume $u = av$ for some $v \in b^*$.
	Let $m = \gcd(|u|_b, n)$.
	Then $v$ induces $m$ cycles of size $n / m$
	on the states. As $a$ only moves the state $s$, for $m > 1$,
	we get $\Syn(\mathcal A_c) \cap u^* = \emptyset$.
    Also, since $d + |u|_b \equiv 0 \pmod{n}$
    implies $d + |u|_b \equiv 0 \pmod{m}$
    and so $d \equiv 0 \pmod{m}$,
    for $m > 1$, we have $d + |u|_b \not\equiv 0 \pmod{n}$.    
    So, the claim follows
    when it holds for $m = 1$,
    which we assume in the following.
	Then $d \equiv j\cdot |u|_b \pmod{n}$ for some $0 < j < n$.
	%	which is coprime to $n$ by \cref{prelim:eqn:ab_divisor}.
	So,~$a$ induces $\SI{s}{\delta(s, v^j)}$
	and $v$ induces $(s ~ \delta(s, v) ~ \ldots ~ \delta(s, v^{n-1}))$
	as transformations on the states.
	By \cref{prelim:eqn:SIcycle2},
	the largest set permuted 
	by $av$ is $\{ \delta(s, v^{j+1}), \delta(s, v^{j+2}), \ldots, \delta(s, v^n)\}$.
	Hence,
	$\Syn(\mathcal A_c) \cap u^* \ne \emptyset$
	if and only if $\delta(s, v^{j+1}) = \delta(s, v^n) = s$,
	which is equivalent to $(j+1)|u|_b \equiv 0 \pmod{n}$.
	As $(j+1)|u|_b \equiv d + |u|_b \pmod{n}$,
	$
	\Syn(\mathcal A_c) \cap u^* \ne \emptyset
	$
	if and only if
	$
	d + |u|_b \equiv 0 \pmod{n}.
	$	
\end{proof}

The condition $d + |u|_b \equiv 0 \pmod{n}$
implies that $|u|_b$ and~$n$ are coprime.
For certain words~$u$, this is, in fact, a necessary condition.
For example, the word $bab^{m - 1}$
for a divisor~$m > 1$ of~$n$ cannot synchronize~$\mathcal A_c$
by \cref{constrsync:subclasses:prop:SI:permutesset-new}.

\begin{corollary}
	\label{constrsync:subclasses:SI:cor:uAc}
	Let $u \in b^*ab^*$ and $i \ge 1$ coprime to $n$. Then $(u+b^i)^* \cap \Syn(\mathcal A_c) \ne \emptyset$.	
\end{corollary}
\begin{proof}
	By coprimality of $i$ and $n$ there exists $j \ge 0$
	such that $|u|_b + j \cdot i + d \equiv 0 \pmod{n}$.
	Then set $v = ub^{j\cdot i}$
	and apply \PROP~\ref{constrsync:subclasses:SI:prop:uAc}.
\end{proof}

\begin{corollary}
	\label{constrsync:subclasses:SI:abb} 	\label{constrsync:SI:base-cases}
	Let $\semiautomat{A_t}{Q}{\{a,b\}}{\delta}$ be an $n$-state SI-automaton
	with %\footnote{A two-state automaton with a simple idempotent letter is always synchronizing.}
	$n \ge 4$,
	simple idempotent letter $a$
	and permutational letter $b$. 	Then the following are equivalent.
	\begin{enumerate}
	\item \label{constrsync:SI:item:abb} $\Syn(\mathcal A) \cap (a+bb)^* \ne \emptyset$,
		
	\item \label{constrsync:SI:item:bab} $\Syn(\mathcal A) \cap (ba^*b)^* \ne \emptyset$,
		
	\item \label{constrsync:SI:item:AcAt} $\mathcal A = \mathcal A_t$ and $n$ is even or $\mathcal A = \mathcal A_c$ and $n$ is odd.	
	\end{enumerate}
\end{corollary}
\begin{proof}
$\ref{constrsync:SI:item:abb} \Leftrightarrow \ref{constrsync:SI:item:bab}$:
As $\fact{(a+bb)^*} = \fact{(ba^*b)^*}$,
this follows from~\cref{constr_sync:cor:factors_sets_included}.

$\ref{constrsync:SI:item:AcAt} \Rightarrow \ref{constrsync:SI:item:bab}$:
	For $\mathcal A = \mathcal A_t$
	and $n$ even, by~\cref{constrsync:subclasses:SI:prop:u-coprime-bs},	
	we have $\Syn(\mathcal A_t) \cap (ba^*b)^* \ne \emptyset$.
	For $\mathcal A = \mathcal A_c$
	and $n$ odd, by~\cref{constrsync:subclasses:SI:cor:uAc},
	we have $\Syn(\mathcal A_c) \cap  (a+bb)^* \ne \emptyset$
	and the claim follows as~\ref{constrsync:SI:item:abb} 
	and~\ref{constrsync:SI:item:bab} are equivalent.
	
$\ref{constrsync:SI:item:bab} \Rightarrow \ref{constrsync:SI:item:AcAt}$:
	By~\cref{prop:structure_simp_idem_binary},
	we have $\mathcal A = \mathcal A_t$
	or $\mathcal A = \mathcal A_c$.
	By~\cref{constrsync:subclasses:cor:SI:permutesset-new} 
	we find that in the former case $n$ must be even,
	and in the latter case $n$ must be odd.		
\end{proof}

Lastly, we collect various constraint languages
which give a problem equivalent to the unconstrained
problem for a specific choice of a simple idempotent
letter and a permutational letter.

\begin{lemma}\label{constrsync:subclasses:SI:equiv-unconstrained}
	Let $\semiautomat{A}{Q}{\{a,b\}}{\delta}$ be an $n$-state SI-automaton
	with simple idempotent letter~$a$
	and permutational letter $b$.
	Then the following are equivalent:
	\begin{enumerate}
		\item \label{constrsync:subclasses:SI:equiv-unconstrained:1} $\Syn(\mathcal A) \ne \emptyset$,
		
		\item \label{constrsync:subclasses:SI:equiv-unconstrained:2} $\Syn(\mathcal A) \cap (bb^*a)^* \ne \emptyset$,
		
		\item \label{constrsync:subclasses:SI:equiv-unconstrained:3} $\Syn(\mathcal A) \cap (ab^*a)^* \ne \emptyset$,
		
		\item \label{constrsync:subclasses:SI:equiv-unconstrained:4} $\Syn(\mathcal A)  \cap (b+aa)^* \ne \emptyset$,
		
		\item \label{constrsync:subclasses:SI:equiv-unconstrained:5} $\Syn(\mathcal A)  \cap (b+ab)^* \ne \emptyset$,
		
		\item \label{constrsync:subclasses:SI:equiv-unconstrained:6} $\Syn(\mathcal A) \cap (bb+ab)^* \ne \emptyset$,
		
		\item \label{constrsync:subclasses:SI:equiv-unconstrained:7} $\Syn(\mathcal A) \cap (bb+ba)^* \ne \emptyset$.
	\end{enumerate}	

\end{lemma}
\begin{proof}
$\ref{constrsync:subclasses:SI:equiv-unconstrained:1} 
\Leftrightarrow \ref{constrsync:subclasses:SI:equiv-unconstrained:2} 
\Leftrightarrow \ref{constrsync:subclasses:SI:equiv-unconstrained:3} 
\Leftrightarrow \ref{constrsync:subclasses:SI:equiv-unconstrained:4} 
\Leftrightarrow \ref{constrsync:subclasses:SI:equiv-unconstrained:5}$: As $a$ is a simple idempotent
	we can replace, for any $i > 0$, each~$a$ by~$a^i$.
	This gives the following equivalences:
	\begin{multline*}
		\Syn(\mathcal A) \ne \emptyset 
		\Leftrightarrow (bb^*a)^* \cap \Syn(\mathcal A) \ne \emptyset 
		\Leftrightarrow (ab^*a)^* \cap \Syn(\mathcal A) \ne \emptyset \\
		\Leftrightarrow (b+aa)^* \cap \Syn(\mathcal A) \ne \emptyset  
		\Leftrightarrow (b+ab)^* \cap \Syn(\mathcal A) \ne \emptyset.
	\end{multline*}

$\ref{constrsync:subclasses:SI:equiv-unconstrained:1} \Leftrightarrow 
\ref{constrsync:subclasses:SI:equiv-unconstrained:6}$:
	Suppose $\Syn(\mathcal A) \ne \emptyset$.
By~\cref{prop:structure_simp_idem_binary} 
we have that $\mathcal A$ equals~$\mathcal A_t$
or~$\mathcal A_c$. Set
\[
w = \begin{cases}
	(ab)^{n-1}         & \text{if } \mathcal A = \mathcal A_t, \\
	(ab^{n-d})^{n-2}ab & \text{if } \mathcal A = \mathcal A_c \text{ and $n$ is even, } \\
	(a(b(b^{n+1})^{n-d-1})^{n-2}ab & \text{if } \mathcal A = \mathcal A_c \text{ and $n$ is odd. } \\ 
\end{cases}
\]
For odd $n$ and $\mathcal A = \mathcal A_c$, $b^{n+1} \in (bb)^*$ induces the same transformation as $b$
on the states. Then, by \cref{prelim:prop:binary-SI-sync-words},
in each case, we have $w \in \Syn(\mathcal A) \cap (bb+ab)^*$.
This yields that $\Syn(\mathcal A) \ne \emptyset$
if and only if $\Syn(\mathcal A) \cap (bb+ab)^* \ne \emptyset$.

	
$\ref{constrsync:subclasses:SI:equiv-unconstrained:6} \Leftrightarrow 
\ref{constrsync:subclasses:SI:equiv-unconstrained:7}$:
As $\fact{(bb+ab)^*} = \fact{(bb+ba)^*}$,
this follows from~\cref{constr_sync:cor:factors_sets_included}.
\end{proof}

% oder einfach annehmen, keine zwei a's hintereinander?
% also wenn a \in Sigma_{1,2}, dann a nicht in Sigma22 oder Sigma23? ne, nur Sigma22

\subsection{Bounded Constraint Languages}
\label{constrsync:input:SI:sec:bounded} 

Here, \cref{constrsync:subclasses:SI:prop:rho} 
states a lower bounds in terms of $\rho \colon \Sigma^* \to \mathbb N_0$
for the size of subsets 
reachable from a given set (recall that $\rho \colon \Sigma^* \to \mathbb N_0$ from Section~\ref{constr_sync:sec:growth} 
measures the number of symbol switches).
Then, in \cref{constrsync:subclasses:SI:cor:rho}, 
we show that we can decide in constant time whether
a given SI-automaton has a synchronizing
word in a fixed letter-bounded language.
Lastly, in \cref{constrsync:prop:SI:boundedLcase},
we show that deciding whether
a given SI-automaton has a synchronizing
word in a regular language in $u_1^* \cdots u_k^*$
with $\rho(u_j) \le 1$ for each $j \in \{1,\ldots,k\}$
is in \LOGSPACE.

\begin{proposition}\label{constrsync:subclasses:SI:prop:rho}
	%\label{constr_sync:prop:SI-automata_equal_letter_runs}
	Let $\semiautomat{A}{Q}{\Sigma}{\delta}$ be an SI-automaton, $S \subseteq Q$
	and $u \in \Sigma^*$. Then $|S| - \rho(u) - 1 \le |\applytrans{\delta}{S}{u}|$.	
	In particular, $|\applytrans{\delta}{S}{u}| = 1$
	implies $\rho(u) \ge |S| - 2$. 
\end{proposition}
\begin{proof}
	First, let $i \ge 0$ be maximal
	such that $u = vx^i$ for some $v \in \Sigma^*$ and $x \in \Sigma$.
	%We can assume $u \ne \varepsilon$, so that $i > 0$.
	Since $x$ is a permutational letter or a simple idempotent,
	we have
	\begin{equation} \label{eqn:SIrho}
	|\delta(S, v)| - 1 \le |\delta(S, u)|.
    \end{equation}
    
    Next, to prove the statement, we do induction on $\rho(u)$.
    For $\rho(u) = 0$, we have $v = \varepsilon$
    and \cref{eqn:SIrho} gives $|S| - 1 \le |\delta(S, u)|$. % the claim.
	Next, suppose $\rho(u) > 0$.
	Then $v$ ends with a letter different from $x$,
	so $\rho(u) = \rho(v) + 1$.
    Now, we have the following:
	\begin{align*}
		|S| - \rho(u) - 1 & = |S| - \rho(v) - 1 - 1 \tag{$\rho(u) = \rho(v) + 1$} \\
		                  & \le |\delta(S, v)| - 1 \tag{induction hypothesis}\\
		                  & \le |\delta(S, u)|.   \tag{Eqn.~\eqref{eqn:SIrho}}
	\end{align*}
	Thus, by induction, we obtain the claim.
\end{proof}

\begin{corollary}\label{constrsync:subclasses:SI:cor:rho}
	Let $L \subseteq \Sigma^*$ be a fixed letter-bounded language.
	
	\begin{enumerate}
	\item \label{SI:cor:rho:item:1} There exist only finitely many SI-automata $\mathcal A$
	with $\Syn(\mathcal A) \cap L \ne \emptyset$.%	having a synchronizing word in~$L$.	
	
	\item \label{item:SI:cor:rho:constant} For a given SI-automaton we can decide in constant time
	whether it has a synchronizing word in $L$.
	\end{enumerate}	
\end{corollary}
\begin{proof}
\ref{SI:cor:rho:item:1} By~\cref{constr_sync:prop:letter_bounded_iff_equal_letter_runs_bounded}
we have $\applyfunc{u}{\rho} \le N$ for some $N \ge 0$
and all $u \in L$.
Then, \cref{constrsync:subclasses:SI:prop:rho}
gives that for an SI-automaton $\semiautomat{A}{Q}{\Sigma}{\delta}$
we have $|\delta(Q, u)| \ge |Q| - N - 1$
so that for a synchronizing word $u \in \Sigma^*$
we find $N + 2 \ge |Q|$.

\ref{item:SI:cor:rho:constant} By~\ref{SI:cor:rho:item:1}, to test whether a given SI-automaton
has a synchronizing word in~$L$
we only have to compare it against finitely many automata, which can be done
in constant time.		
\end{proof}

\begin{proposition} \label{constrsync:prop:SI:boundedLcase}
	Let $L \subseteq u_1^* \cdots u_k^*$
	be regular with $u_1, \ldots, u_k \in a^* b^* \cup b^* a^*$.
	Then deciding whether an SI-automaton $\semiautomat{A}{Q}{\{a,b\}}{\delta}$ 
	has a synchronizing word in~$L$	is in~\LOGSPACE.
\end{proposition}
\begin{proof}
	Let $n$ be the number of states of~$\mathcal A$.
	We can assume $n \ge 1$ and that~$\mathcal A$ has 
	a simple idempotent letter and a permutational letter.	
	Since~$L$ is regular, it has a finite number~$m$ of Myhill-Nerode equivalence
	classes by \cref{prelim:min_aut:terminal_object}.
	
	Let $u_1^{m_1} \cdots u_k^{m_k} \in L$ be a shortest
	synchronizing word for $\mathcal A$.	
	Assume we have $m_j \ge nm$ for some $j \in \{1,\ldots,k\}$.
	By finiteness, we have
	\[
	u_1^{m_1} \cdots u_{j-1}^{m_{j-1}} u_j^{i_1} \equiv_{L} 
	u_1^{m_1} \cdots u_{j-1}^{m_{j-1}} u_j^{i_1 + p_1}
	\]
	for some $0 \le j_1 < m$ and $0 < p_1 < m$.
	Now, $u_j$ either is a permutational word, a simple idempotent,
	or the composition of a permutational word with a simple idempotent.
	With \cref{constrsync:prop:concat_cycle_SI} 
	we have $\delta(q, u_j^{i_2}) = \delta(q, u_j^{i_2 + p_2})$
	for all $q \in Q$ for some $0 \le i_2 < n$
	and $0 < p_2 < n$.
	Since 
	\[
	nm - p_1 p_2 = (n- p_2) m + p_2(m - p_1) = n ( m - p_1 ) - (n - p_2) p_1
	\]
	we have $m_j - p_1 p_2 \ge nm - p_1 p_2 \ge \max\{ j_1, j_2 \}$.
	Thus, we find
	\[
	u_1^{m_1} \cdots u_j^{m_j - p_1 p_2} \cdots u_k^{m_k} \equiv_L 
	u_1^{m_1} \cdots u_j^{m_j} \cdots u_k^{m_k}
	\]
	and $\delta(q, u_j^{m_j - p_1 p_2}) = \delta(q, u_j^{m_j})$
	for all $q \in Q$.
	This yields that $u_1^{m_1} \cdots u_j^{m_j - p_1 p_2} \cdots u_k^{m_k} \in L$
	and that this word synchronizes~$\mathcal A$,
	which contradicts the minimal choice of the original word.
	Thus, we have $m_j < nm$ for all $j \in \{1,\ldots,k\}$.
		
	Hence, there exists a synchronizing word in $L$ if and only if there exists one of the form $u_1^{m_1} \cdots u_k^{m_k}$ in $L$
	with $0 \le m_1, \ldots, m_k < nm$, and so we only have to search among these words.
	By utilizing~$k$ counters each 
	using $\log nm$ space,
	we iterate through all tuples $(m_1, \ldots, m_k)$
	representing the words $u_1^{m_1} \cdots u_k^{m_k}$
	with $0 \le m_1, \ldots, m_k < nm$.
	In each iteration, we test whether  $u_1^{m_1} \cdots u_k^{m_k}$
	synchronizes~$\mathcal A$ and is in $L$.
	The word $u_1^{m_1} \cdots u_k^{m_k}$
	can be tested to be synchronizing for $\mathcal A$ in deterministic log-space
	by iterating through $Q = \{q_1, \ldots, q_n\}$,
	compute $\delta(q_i, u_1^{m_1} \cdots u_k^{m_k})$ on-the-fly
	and test, for $i > 1$, whether $q_{i-1}$ equals $q_i$.
	This only requires two additional memory cells and three counters for 
	the state $q_i$, $i \in \{1,\ldots,n\}$, for the index $j \in \{1,\ldots,k\}$
	and for $u_j^{i}$ with $0 \le i \le m_j$.
	Checking whether $u_1^{m_1} \cdots u_k^{m_k} \in L$
	with input $(m_1, \ldots, m_k)$ can be done
	by simulating an automaton accepting~$L$ from the start state while using
	counters as before.
\end{proof} 

\subsection{Binary SI-Automata and Small Constraint Automata}
\label{constrsync:input:SI:sec:constraintaut}

Here, first, in \cref{constrsync:prop:SI-binary-sync-in-L-hard} 
we show \LOGSPACE-hardness of deciding whether
a given binary SI-automaton has a synchronizing word in 
a certain constraint language. Recall
that by \cref{syncmax:SI-aut:prop:sync-L-complete}
the unconstrained synchronization problem is \LOGSPACE-complete
on binary SI-automaton as input.
In \cref{constrsync:input:SI:lem:threestates} 
we show that for binary $3$-state constraint PDFAs $\mathcal B$
with two strongly connected components,
deciding
whether a given SI-automaton has a synchronizing
word in $L(\mathcal B)$ is in~\LOGSPACE.
Lastly, \Cref{constrsync:input:SI:thm:atmostthreestates} 
states that the constrained synchronization problem
for binary constraint PDFAs with at most three states and SI-automata
as input is in \NL. 

\begin{proposition} \label{constrsync:prop:SI-binary-sync-in-L-hard}
	Let $L \subseteq \{a,b\}^*$ 
	be such that $v u^* w \subseteq L$ for some $v,w \in \Sigma^*$
	and $u \in \Sigma^*a\Sigma^* \cap \Sigma^*b\Sigma^*$.
	Then it is \LOGSPACE-hard w.r.t.~$\leACZEROm$-reductions
	to test whether a given input SI-automaton
	has a synchronizing word in $L$.	
\end{proposition}
\begin{proof}
	We give a reduction from \SCP, which is \LOGSPACE-complete
	by \cref{prelim:thm:complexities}.
    Let $f \colon [m] \to [m]$
	be a permutation. 
	By modifying $f$ by  
	adding a new point $n$
	and mapping $m-1$ to $m$
	and $m$ to $f(m-1)$,
	we obtain a permutation on $[m+1]$
	that cyclically permutes its domain
	if and only if $f$ cyclically permutes $[m]$.
    So, assume that~$m$ and~$|u|_b$
	are coprime.
	Then, let $\semiautomat{A}{Q}{\{a,b\}}{\delta}$
	with $Q = [m] \cup \{t\}$ where~$a$ induces a simple idempotent mapping $0$ to $t$, $\delta(t, b) = t$
	and $\delta(q, b) = f(q)$ for $q \in [m]$.
		
	For the input $1^m$, since coprimality can be tested in \LOGSPACE 
	for unary inputs by \cref{prelim:prop:coprime_in_L}, a Turing machine can determine how many points must be added 
	to the input $f$ so that the total number of points is coprime to $|u|_b$, 
	and can compute an $\AC^0$ circuit realizing this modification of $f$ 
	and computing~$\mathcal{A}$. Hence, we obtain an~$\leACZEROm$-reduction.	
		
%We give a reduction from \SCP. Let $f \colon [n] \to [n]$
%be a permutation.
%Let $0 < m \le |u|_b+1$
%such that $n + m \equiv 1 \pmod{|u|_b}$. % uniform TM, rest |u|_b bestimmen, dann draufaddiern
%Set $\semiautomat{A}{Q}{\{a,b\}}{\delta}$
%with $Q = [n] \cup \{q_1, \ldots, q_m, t\}$ where~$a$ induces a simple idempotent mapping $0$ to $t$
%and
%
%\[
%\delta(q, b) = \begin{cases}
%	q_1     & \text{if } q = 0, \\
%	q_{i+1}	& \text{if } q = q_i \text{ for } i \in \{1,\ldots,m-1\}, \\
%	f(0)    & \text{if } q = q_m, \\
%	f(q)	& \text{if } q \in \{1,\ldots,n-1\}, \\
%	t       & \text{if } q = t.
%\end{cases}
%\]
%
%A Turing machine on input $1^n$
%can compute $m$ by computing the remainder~$r$
%of $n$ when divided by $|u|_b$
%and setting $m = |u|_b - r + 1$.
%Then it can output an~$\AC^0$ circuit computing $\mathcal A$.
%So, we obtain an $\leACZEROm$-reduction.

	Now, $f$ cyclically permutes $[m]$
	if and only if $\mathcal A$ is isomorphic 
	to the automaton from \cref{prelim:fig:binarySI-sync-trap}.	
	So, if~$f$ cyclically permutes $[m]$, by \cref{constrsync:subclasses:SI:prop:u-coprime-bs},
	as $|Q|-1$ and $|u|_b$ are coprime
	we have $u^i \in \Syn(\mathcal A)$
	for some $i \ge 1$.
	Then, $vu^iw \in \Syn(\mathcal A)$.
	Conversely, if $L \cap \Syn(\mathcal A) \ne \emptyset$,
	then, as $\mathcal A$ has a trap state, it 
	must be isomorphic to the automaton from \cref{prelim:fig:binarySI-sync-trap}.
	So, $f$ cyclically permutes $[m]$.
	Hence, we have a valid reduction.	
\end{proof} 


\begin{lemma} \label{constrinp:lem:three-states-two-SCC}
	Let $\automat{B}{\{1,2,3\}}{\{a,b\}}{\mu}{1}{\{3\}}$
	be a trim $3$-state constraint automaton
	with two strongly connected components.
	Then $L(\mathcal B)$
	is letter-bounded
	or there exist $v,u,w \in \Sigma^*$
	with $vu^*w \subseteq L(\mathcal B)$
	and $|u|_a, |u|_b > 0$.
\end{lemma}
\begin{proof}
	Let $p \in \{1,2,3\}$
	be a state in the strongly connected component
	of size two and~$q$
	be the single state in the other strongly connected component.
	By \cref{constr_sync:thm:regular_prefix_codes_and_automata}
	we have $L(\mathcal B_{p,\{p\}}) = C_1^*$
	and $L(\mathcal B_{q,\{q\}}) = C_2^*$
	where $C_1$ is a prefix code (the case $C_1 = \emptyset$ is not possible
	since the strongly connected component has size two)
	and $C_2 = \emptyset$ or it is also a prefix code.
	Now, $L(\mathcal B)$
	is letter-bounded if and only if $C_1 \subseteq x^*$
	and $C_2 \subseteq y^*$ for $x,y \in \Sigma$.
	Otherwise, $C_1^*$ or $C_2^*$
	contains a word $u$
	with $|u|_a, |u|_b > 0$.
	Without loss of generality, suppose $u \in C_1^*$.
	We have $u^* \subseteq C_1^*$
	and since~$\mathcal B$
	is trim,
	we have $\mu(1, v) = p$
	and $\mu(p, w) = 3$
	for some $v,w \in \{a,b\}^*$.
	This gives $vu^*w \subseteq L(\mathcal B)$.
\end{proof}

\begin{lemma}\label{constrsync:input:SI:lem:threestates}
	Let $\automat{B}{\{1,2,3\}}{\{a,b\}}{\mu}{1}{\{3\}}$
	be a trim $3$-state constraint automaton
	with two strongly connected components.
	Then the problem to decide
	whether a given SI-automaton over $\{a,b\}$
	has a synchronizing word in $L(\mathcal B)$ is in \LOGSPACE.
\end{lemma}
\begin{proof}
	Let $\semiautomat{A}{Q}{\{a,b\}}{\delta}$
	be an $n$-state input SI-automaton.  
	As finitely many cases can be hard-coded into a decision algorithm, we can 
	assume $n \ge 24$.  	
	Then, by \cref{prop:structure_simp_idem_binary}, 
	a necessary condition for $\mathcal A$
	having a synchronizing word in~$L$
	is that it equals, up to renaming of the letters,
	$\mathcal A_c$
	or $\mathcal A_t$.
	By \cref{constrsync:prop:SI-binary-sync-in-L}
	we can test in deterministic log-space whether this is the case
    and compute the respective parameters.
    So, we can assume 
    $\mathcal A$ equals $\mathcal A_c$ or $\mathcal A_t$,
    where $a$ is a simple idempotent and $b$
    the permutational letter.
    Further, assume 
    $\fact{L(\mathcal B}) \ne \{a,b\}^*$,
    as otherwise, by~\cref{constr_sync:cor:factors_sets_included},
    when $\mathcal A$ equals $\mathcal A_c$
    or $\mathcal A_t$ it is $\Syn(\mathcal A) \cap L(\mathcal B) \ne \emptyset$.    
    %the problem is solved after testing $\Syn(\mathcal A) \ne \emptyset$.    
	Next, set $\Sigma_{i,j} = \{ x \in \Sigma \mid \mu(i, x) = j \}$
	for $i,j \in \{1,2,3\}$.
	Let $S \subseteq \{1,2,3\}$
	be the strongly connected component with two states.
	Then 
	\begin{equation}\label{constrsync:SI:eqn:threestates}
		L(\mathcal B) = \begin{cases}
			( \Sigma_{1,1}^* \Sigma_{1,2} \Sigma_{2,2}^* \Sigma_{2,3} \cup \Sigma_{1,1}^* \Sigma_{1,3} ) L(\mathcal B_{3,\{3\}}) & \text{if } S = \{2,3\}, \\
			L(\mathcal B_{1,\{1\}}) ( \Sigma_{1,3}\Sigma_{3,3}^* \cup \Sigma_{1,2}\Sigma_{2,2}^* \Sigma_{2,3} \Sigma_{3,3}^*) & \text{if } S = \{1,2\}.
		\end{cases}
	\end{equation}
    So, $L(\mathcal B) = ULV$ for letter-bounded $U, V \subseteq \Sigma^*$
    with $\rho(u) \le 4$ for $u \in U \cup V$
    and~$L$ equals $L(\mathcal B_{3,\{3\}})$ for $S = \{2,3\}$
    or $L(\mathcal B_{1,\{1\}})$ for $S = \{1,2\}$.    
    As $\fact{L(\mathcal B}) \ne \{a,b\}^*$, we have
    $|\Sigma_{i,i}| \le 1$ for all $i \in \{1,2,3\}$
    and only two or three transitions
    are defined among the states in $S$.
    This yields that there are only~$13$
    possible cases for~$L$, which we handle next.  
    Observe that, as $L \subseteq \fact{L(\mathcal B)}$,
    by \cref{constr_sync:thm:L_in_FactLbar},
    $\Syn(\mathcal A) \cap L \ne \emptyset$
    implies $\Syn(\mathcal A) \cap L(\mathcal B) \ne \emptyset$.     
    
    \medskip
    \noindent\underline{$(ab)^*, (ba)^*, (a+ba)^*, (aa^*b)^*, (ab+aa)^*$}:
       As $a$ is a simple idempotent,
       we can replace $(a+ba)^*$, $(aa^*b)^*$
       or $(ab+aa)^*$ by the bounded languages
       $(a+\varepsilon)(ba)^*$, $(ab)^*$
       or $(ab)^*(a+\varepsilon)$, respectively.
       So, for each language, 
       we can assume that we have $L(\mathcal B) = ULV \subseteq u_1^* \cdots u_k^*$
       for some words $u_i \in \Sigma^{\le 2}$ with $i \in \{1,\ldots,k\}$.
       Then, by \PROP~\ref{constrsync:prop:SI:boundedLcase},
       we find that the decision problem is in~\LOGSPACE, giving the claim.
    
    \medskip
    \noindent\underline{$(bb^*a)^*, (ab^*a)^*, (b+aa)^*, (b+ab)^*, (bb+ab)^*, (bb+ba)^*$}:
    By \cref{constrsync:subclasses:SI:equiv-unconstrained},
    as it lists these languages, $\Syn(\mathcal A) \ne \emptyset$
    is equivalent to $\Syn(\mathcal A) \cap L \ne \emptyset$. 
    As $\Syn(\mathcal A) \cap L \ne \emptyset$
    implies $\Syn(\mathcal A) \cap L(\mathcal B) \ne \emptyset$,
    we find that $\Syn(\mathcal A) \cap L(\mathcal B) \ne \emptyset$
    is equivalent to $\Syn(\mathcal A) \ne \emptyset$,
    which gives the claim.       

    \medskip
    \noindent\underline{$(ba^*b)^*$, $(a+bb)^*$}:
    First, assume $\mathcal A = \mathcal A_c$ and~$n$
    is even (resp.~$\mathcal A = \mathcal A_t$ and $n$ is odd).
    Let  $u_1 \in U$, $u_2 \in L$ 
    and $u_3 \in V$. 
    By \cref{constrsync:subclasses:cor:SI:permutesset-new},
    there exists $S \subseteq Q$
    with $|S| \ge n/2 \ge 12$ (resp.~$|S| \ge (n-1)/2 \ge 12$)
    such that $\delta(S, u_2) = S$.
    As $\rho(u_1), \rho(u_3) \le 4$, by~\cref{constrsync:subclasses:SI:prop:rho},
    we get
    \[
    |\delta(Q, u_1)| \ge n - 5 \quad\text{ and }\quad 
    |\delta(Q, u_1 u_2 u_3)| \ge |\delta(Q, u_1 u_2)| - 5.
    \]
    The former equation yields $|\delta(Q, u_1) \cap S| \ge 7$.
    As $u_2$ permutes the states in~$S$
    we have $|\delta(Q, u_1u_2) \cap S| = |\delta(Q, u_1) \cap S|$.
    Then the latter equation yields $|\delta(Q, u_1 u_2 u_3)| \ge |\delta(Q, u_1 u_2) \cap S| - 5 \ge 2$.
    So, $\Syn(\mathcal A) \cap L(\mathcal B) = \emptyset$.    
    Otherwise, if $\mathcal A = \mathcal A_c$ and $n$ is odd
    (resp.~$\mathcal A = \mathcal A_t$ and $n$ even),
    then, by \cref{constrsync:subclasses:SI:abb},
    we have $\Syn(\mathcal A) \cap L \ne \emptyset$.
    As written above, this implies $\Syn(\mathcal A) \cap L(\mathcal B) \ne \emptyset$.
    Now, whether $n$ is odd or even can be decided 
    in deterministic log-space.    
\end{proof}

\begin{proposition} \label{constrsync:input:SI:prop:threestates-not-SCC}
	Let $\mathcal B = (P, \{a,b\}, \mu, p_0, F)$
	be a constraint automaton with $|P| \le 3$
	that is not strongly connected.
	For an SI-automaton $\semiautomat{A}{Q}{\{a,b\}}{\delta}$
	the problem to decide whether it has a synchronizing
	word is decidable in constant time if $L(\mathcal B)$
	is letter-bounded and it is \LOGSPACE-complete
	otherwise.
\end{proposition}
\begin{proof}
	For $|P| \le 2$
	or if $\mathcal B$ has precisely three strongly connected components and three states,
	then it is easy to see that either $L(\mathcal B)$ is letter-bounded
	or $\factor{L(\mathcal B)} = \{a,b\}^*$ (in fact, in the latter case, 
	we have $L(\mathcal B_{p,\{p\}}) = \{a,b\}^*$
	for some $p \in P$ that can assumed to be final,
	which gives $\factor{L(\mathcal B)} = \{a,b\}^*$,
	or for all $p \in P$ we have $L(\mathcal B_{p,\{p\}}) \subseteq x^*$
	for some $x \in \Sigma$,
	and so the accepted language is letter-bounded by \crefitem{constr_sync:prop:letterboundedreg}{constr_sync:prop:letterboundedreg:item:uvw}).	
	In the latter case,
	by \cref{constr_sync:cor:factors_sets_included} 
	the problem is the same as deciding whether the automaton
	is synchronizing, and so \LOGSPACE-complete
	by \cref{syncmax:SI-aut:prop:sync-L-complete}.
	In the former case, the problem is solvable in constant time
	by \crefitem{constrsync:subclasses:SI:cor:rho}{item:SI:cor:rho:constant}.

	Lastly, if $|P| = 3$ and~$\mathcal B$
	has two strongly connected components,
	then we can assume $F = \{3\}$
	by \cref{constr_sync:prop:final_states,gen_results:cor:factors_etc_complexity}
	and since~$\mathcal B$ is trim.
	Then \cref{constrinp:lem:three-states-two-SCC,constrsync:input:SI:lem:threestates,constrsync:prop:SI-binary-sync-in-L-hard,constrsync:subclasses:SI:cor:rho} give the claim.	
\end{proof}

\Cref{constrsync:input:SI:prop:threestates-not-SCC} shows that for several constraint languages
that yield \PSPACE-complete problems in general,
for SI-automata as input the complexity drops to containment in \LOGSPACE-completeness.

\begin{theorem}\label{constrsync:input:SI:thm:atmostthreestates}
		Let $\mathcal B = (P, \{a,b\}, \mu, p_0, F)$
		be a constraint automaton with $|P| \le 3$.
		For an $n$-state SI-automaton $\semiautomat{A}{Q}{\{a,b\}}{\delta}$
		we can decide in \NL\  whether it has a synchronizing word in $L(\mathcal B)$.
\end{theorem}
\begin{proof}
	If~$\mathcal B$ is strongly connected,
	then the claim follows from \cref{constrsync:cor:SCC}.
	Otherwise, it follows from \cref{constrsync:input:SI:prop:threestates-not-SCC}.
\end{proof}